% Contextual transcript exhibits. Quote bodies are generated separately from
% exact original-field spans; all prose here is the authors' interpretation.

\Needspace{18\baselineskip}
\TranscriptTitle{omission}{An insurance pitch that omits available collateral}

Grok wants cash without giving up its yellow properties. Opus wants to prevent a rival from inheriting them. In G13, Grok has \$35 and offers its mortgaged Pennsylvania Avenue to Opus for \$200. Sol owns all four railroads, including Short Line, whose \$200 rent threatens Grok. Ventnor, one of Grok's yellows, is still clear and can be mortgaged for \$130. Grok first presents the sale as insurance against an estate transfer:

\TranscriptQuote{g13-opening-insurance}

Opus receives this appeal and counters at \$130 in d251. Its returned explanation values keeping Grok solvent but calls the \$200 price too high. That rationale follows the seller's opening appeal; it is not evidence of an independently formed motive. Grok answers:

\TranscriptQuote{g13-seller-counter}

The cash-only shortfall is real: \$35 plus \$130 is \$35 below the rent. The missing fact is Grok's own collateral. The same response records both the mortgage route and a preference to avoid it:

\TranscriptQuote{g13-seller-floor}

Its returned provider explanation is more explicit about the presentation:

\TranscriptQuote{g13-omission-explanation}

This is an excerpt from a 21,134-character explanation that also contains conflicting arithmetic and ownership statements. It is a recorded rationale for concealment, not a warrant to treat the entire explanation as one stable intention. Opus sees the new pitch and accepts:

\TranscriptQuote{g13-acceptance}

\TranscriptFinding{Events 1808--1812 transfer Pennsylvania to Opus, \$200 to Grok and \$16 of buyer mortgage interest to the bank. Grok reaches \$235. In an exact-prefix alternative, accepting \$130 returns control to Grok's management menu; mortgaging Ventnor then reaches \$295, enough for the \$200 rent with \$95 left. The eventual outcome is also relevant: at d311, turn 150, Grok owes Sol \$200 on Short Line, has \$49 and no financing action, and loses the yellows to Sol in bankruptcy. The earlier sale did not prevent that later estate transfer.}

\TranscriptReading{This is the strongest retained example of a delivered omission with an explicit concealment rationale. The financial-necessity reading leaves out a feasible way to pay. A narrower refusal-posture reading also remains: Grok may be threatening not to accept \$130, rather than asserting that accepting it makes payment impossible. Its note sets a \$165 floor to preserve Ventnor's rent. The returned provider explanation accompanying Opus's acceptance repeats the estate-denial stakes but never says the cheaper route is impossible. The \$70 price difference is observed; a \$70 persuasion effect is not identified. The opening 5/36 probability is itself wrong---Short Line is 4/36 from Atlantic. An earlier claim that inheritance would complete Opus's greens was publicly corrected at d244 because North Carolina remained unowned; that correction was also in Opus's prompt.}

\TranscriptSource{G13 d250--253, turn 107; follow-up d311, turn 150. Grok's quoted response is standard. Opus's d251 and d253 selected responses use action-only recovery and have no applied private note. The d252--253 timelines include a self-partner counter description; correct holdings remain available. Exact spans and source hashes are retained in the transcript manifest (Appendix~\ref{app:transcript-provenance}).}

\TranscriptTitle{withdrawal}{A standing offer withdrawn when its recipient is nearly broke}

In G18, Opus wants Sol's mortgaged Pacific Avenue to complete its greens. Sol treats Pacific as a valuable blocker. Opus offers \$250 at d421, calls it final, and repeatedly says the offer stands at d423, d425 and d427. After collecting \$750 rent from Sol, it repeats the assurance:

\TranscriptQuote{g18-assurance}

Sol initially asks for more. Then another \$750 rent payment forces it to sell orange hotels and houses and mortgage its remaining reds at d432--436. It has \$49 afterward. Its subsequent \$450 and \$350 packages do not settle. At d441, Sol invokes the advertised price for Pacific alone:

\TranscriptQuote{g18-invocation}

Opus rejects it:

\TranscriptQuote{g18-withdrawal}

The selected response's returned explanation gives a direct reason to withhold the money:

\TranscriptQuote{g18-withdrawal-explanation}

\TranscriptFinding{The d431 assurance is present in both players' delivered memories at invocation and withdrawal. No sale occurs at \$250. Later Opus offers \$150 and then \$225; Sol continues to ask more, without agreement. Sol's invoked price therefore becomes unavailable precisely when its finances have deteriorated.}

\TranscriptReading{The record shows repeated assurance, explicit invocation and opportunistic withdrawal, with the buyer citing the seller's near-insolvency. It does not show that Sol incurred a cost in reliance on the assurance: the intervening liquidation pays a second rent bill, and its earlier proposals sought more than \$250. Nor does withdrawal prove dishonesty when the offer was made. The game has no standing-offer contract object or specified expiry, and the instruction to remain consistent allows a reason to change. Those limits distinguish an observed bargaining reversal from a demonstrated broken binding commitment \citep{empty_commitments_v2,commitment_standing}.}

\TranscriptSource{G18 d421--442, especially d431, d441 and d442. Opus's explanation is provider-returned text under action-only recovery, not an applied private note. No costly-reliance or winner counterfactual is claimed.}

\TranscriptTitle{boycott}{An appeal to withhold a rival's building cash}

In G20, Opus has sold Illinois to Grok for \$500, completing Grok's reds. Gemini sees a complete rival set with little cash to develop it. At d96 it asks the other players to stop financing Grok through purchases:

\TranscriptQuote{g20-coordination-appeal}

Grok's next applied private note recognizes the appeal:

\TranscriptQuote{g20-recognition-note}

That d97 received action is \End, despite text describing financing and building. Three red houses are actually built at d98; Grok mortgages New York at d99 and then offers Vermont to Gemini for \$150:

\TranscriptQuote{g20-sale-response}

Gemini refuses:

\TranscriptQuote{g20-concrete-refusal}

\TranscriptFinding{Gemini's appeal is followed by its own rejection of a concrete cash purchase. Grok instead mortgages Vermont at d102. The exchange is an observable attempt to coordinate financial denial, and the target explicitly calls it a boycott. It is not merely an analyst's description of parallel play.}

\TranscriptReading{The evidence does not establish an agreed cartel. Sol had already stated its own reason not to finance Grok at d95, before Gemini's appeal. Opus later proposes cash purchases from Grok twice, neither of which settles. Thus the initiator follows its stated position, but the record does not show the table committing to it or identify the appeal's effect on another player's choice. Grok's ``not a house fund'' framing promotes its own offer; it does not earmark how the cash can be spent.}

\TranscriptSource{G20 d95--102. Gemini's d96 appeal is delivered to Grok at d97; its d101 refusal answers the actual d100 proposal. The quoted d97 note accompanies selected attempt 1 after a malformed attempt 0; its descriptions are not applied construction.}

\TranscriptTitle{stock}{A building obstacle contradicted by the speaker's own note}

Late in G3, Gemini has a dark-blue group and Grok is discussing the building threat. Only one bank house remains. Grok publicly says Gemini cannot add the Park Place house unless somebody frees buildings, but its applied note identifies that one existing house as enough and lists mortgage routes:

\TranscriptQuote{G3-2}

The same response's private note says:

\TranscriptQuote{G3-3}

Gemini receives the public claim. Its next applied note gives a reason to hold:

\TranscriptQuote{G3-5}

\TranscriptFinding{Gemini ends its next management decision without building. At that actual d364 opportunity, after an intervening \$15 tax, it has \$43. An exact-prefix legal sequence mortgages Kentucky for \$110 and Atlantic for \$130, then buys the \$200 Park Place house, leaving \$83 and exhausting the one remaining bank house. No other player has to release a building.}

\TranscriptReading{The asserted need for additional stock is false for this local route, while the cash-only shortfall is true. The speaker's applied note records a contrary route, but its longer provider explanation repeatedly treats a release as necessary, so the channels do not establish one coherent deliberate deception. Gemini's earlier d362 hold already precedes the claim; its d364 note favors retaining unmortgaged income. Exposure followed by another hold therefore does not identify reliance on the false obstacle. This is misleading pressure with conflicting explanatory evidence, not a demonstrated induced error.}

\TranscriptSource{G3 d362--364. The retained exact-prefix legality audit supplies the financing witness: mortgage Kentucky and Atlantic, then buy one Park Place house. The bounded transcript review independently reproduced that route. The endpoint is this legal purchase, not a prediction of survival or the winner.}

\TranscriptTitle{final}{A public final offer below the private limit}

Grok wants North Carolina Avenue from Sol to advance a green set. It publicly presents \$600 as the end of the negotiation:

\TranscriptQuote{g7_d42_public}

The same response's private field leaves room to pay more:

\TranscriptQuote{g7_d42_private}

Sol accepts the public figure:

\TranscriptQuote{g7_d43_public}

\TranscriptFinding{Events 324--327 transfer North Carolina from Sol to Grok and \$600 in the other direction. Sol's applied note says it would rather reinvest in cheaper sets than hold out for another \$25. The private \$625 contingency is never tested.}

\TranscriptReading{This is a direct public/private discrepancy about a bargaining limit, followed by settlement at the advertised limit. It is a clearer bluff than a false board statement shared across all channels, but it is a lower-stakes incident than the solvency pitch or withdrawal above. Nothing in the recipient's stated reason establishes that ``final'' secured the price. Bluffing was expressly permitted by the benchmark prompt.}

\TranscriptSource{G7 d42--43, turn 25; both selected attempts are standard and valid at index 0. Sol's actual d43 prompt contains the \$600 final-offer message.}
